% GENERATED FILE. Edit canonical lessons, then run npm run generate:books.
% canonical-source-hash: fnv1a64:c786223c2818dc41
% canonical-lessons: MR-C471-sahana, MR-C471-catur, MR-C471-buddhiman, MR-C471-samjutdar, MR-C471-akkal

\chapter{Good Sense}
\label{ch:mr-good-sense}
\hlchaptermodality{471}

\begin{chapteropening}
\textbf{By the end of this chapter:} \emph{I can describe a sensible, shrewd, intelligent or understanding person.}

Complete the last lesson of chapter 471: I can describe a sensible, shrewd, intelligent or understanding person.
\end{chapteropening}

\section[\texorpdfstring{\mr{शहाणा} (śahāṇā)}{शहाणा (śahāṇā)}]{\mr{शहाणा} (śahāṇā) — sensible; wise}

\label{lesson:MR-C471-sahana}



\begin{quote}
Before the new one: say the Marathi for innocent, then the Marathi for to dupe.
\end{quote}

\subsection*{You'll want to know: \mr{शहाणा}}
\textbf{\mr{शहाणा}} — \emph{śahāṇā} — \textquotedblleft{}sensible; wise\textquotedblright{}.

\begin{quote}
\mr{शहाणी} \mr{मुलगी} \mr{आईला} \mr{कामात} \mr{मदत} \mr{करते}.
\end{quote}

\begin{quote}
\emph{śahāṇī mulgī āīlā kāmāt madat karte.}
\end{quote}

\begin{quote}
The sensible girl helps her mother with the work.
\end{quote}

Agrees: \mr{शहाणा}, \mr{शहाणी}, \mr{शहाणं}. \mr{शहाणं} \mr{बाळ}, a good baby.

\begin{cousinweb}
\textbf{\mr{शहाणा}} (\emph{śahāṇā}) is a Marathi adjective whose origin is not settled.
\end{cousinweb}

\begin{grammarlens}[title={What you've built}]
One more word for describing what you lived and felt.
\end{grammarlens}

\subsection*{Your turn}
\begin{itemize}
\raggedright
  \item \emph{Say it:} \emph{śahāṇā}
  \item \emph{Say it:} \emph{śahāṇā}, once more
  \item \emph{Say it:} \emph{śahāṇī mulgī āīlā kāmāt madat karte.} — The sensible girl helps her mother with the work.
  \item \emph{Recall:} say the Marathi for innocent, then the Marathi for to dupe, then say \emph{śahāṇā} again
\end{itemize}

\begin{tcolorbox}[breakable,colback=teal!4,colframe=teal!35!black,title={Before you move on}]
What does \mr{शहाणा} mean? (\textquotedblleft{}Sensible; wise\textquotedblright{}.) Say it once more.
\end{tcolorbox}
\section[\texorpdfstring{\mr{चतुर} (catur)}{चतुर (catur)}]{\mr{चतुर} (catur) — shrewd; clever}

\label{lesson:MR-C471-catur}



\begin{quote}
Before the new one: say the Marathi for to dupe, then the Marathi for sensible.
\end{quote}

\subsection*{You'll want to know: \mr{चतुर}}
\textbf{\mr{चतुर}} — \emph{catur} — \textquotedblleft{}shrewd; clever\textquotedblright{}.

\begin{quote}
\mr{चतुर} \mr{कोल्ह्याने} \mr{कावळ्याला} \mr{फसवलं}.
\end{quote}

\begin{quote}
\emph{catur kolhyāne kāvḷyālā phasavlaṁ.}
\end{quote}

\begin{quote}
The clever fox tricked the crow.
\end{quote}

Unchanging; \mr{चतुराई}, shrewdness.

\begin{cousinweb}
\textbf{\mr{चतुर}} (\emph{catur}) is a tatsama, taken from Sanskrit \emph{catura}, \textquotedblleft{}quick, clever\textquotedblright{} (CDIAL 4594).
\end{cousinweb}

\begin{grammarlens}[title={What you've built}]
One more word for describing what you lived and felt.
\end{grammarlens}

\subsection*{Your turn}
\begin{itemize}
\raggedright
  \item \emph{Say it:} \emph{catur}
  \item \emph{Say it:} \emph{catur}, once more
  \item \emph{Say it:} \emph{catur kolhyāne kāvḷyālā phasavlaṁ.} — The clever fox tricked the crow.
  \item \emph{Recall:} say the Marathi for to dupe, then the Marathi for sensible, then say \emph{catur} again
\end{itemize}

\begin{tcolorbox}[breakable,colback=teal!4,colframe=teal!35!black,title={Before you move on}]
What does \mr{चतुर} mean? (\textquotedblleft{}Shrewd; clever\textquotedblright{}.) Say it once more.
\end{tcolorbox}
\section[\texorpdfstring{\mr{बुद्धिमान} (buddhimān)}{बुद्धिमान (buddhimān)}]{\mr{बुद्धिमान} (buddhimān) — intelligent}

\label{lesson:MR-C471-buddhiman}



\begin{quote}
Before the new one: say the Marathi for sensible, then the Marathi for shrewd.
\end{quote}

\subsection*{You'll want to know: \mr{बुद्धिमान}}
\textbf{\mr{बुद्धिमान}} — \emph{buddhimān} — \textquotedblleft{}intelligent\textquotedblright{}.

\begin{quote}
\mr{मीरा} \mr{अभ्यासात} \mr{बुद्धिमान} \mr{आहे}.
\end{quote}

\begin{quote}
\emph{mīrā abhyāsāt buddhimān āhe.}
\end{quote}

\begin{quote}
Meera is bright at her studies.
\end{quote}

\mr{बुद्धी}, intellect, + -\mr{मान}, having. Unchanging.

\begin{cousinweb}
\textbf{\mr{बुद्धिमान}} (\emph{buddhimān}) is a tatsama, taken from Sanskrit \emph{buddhimat} (nominative \emph{buddhimān}), \textquotedblleft{}endowed with understanding, intelligent\textquotedblright{} (Monier-Williams).
\end{cousinweb}

\begin{grammarlens}[title={What you've built}]
One more word for describing what you lived and felt.
\end{grammarlens}

\subsection*{Your turn}
\begin{itemize}
\raggedright
  \item \emph{Say it:} \emph{buddhimān}
  \item \emph{Say it:} \emph{buddhimān}, once more
  \item \emph{Say it:} \emph{mīrā abhyāsāt buddhimān āhe.} — Meera is bright at her studies.
  \item \emph{Recall:} say the Marathi for sensible, then the Marathi for shrewd, then say \emph{buddhimān} again
\end{itemize}

\begin{tcolorbox}[breakable,colback=teal!4,colframe=teal!35!black,title={Before you move on}]
What does \mr{बुद्धिमान} mean? (\textquotedblleft{}Intelligent\textquotedblright{}.) Say it once more.
\end{tcolorbox}
\section[\texorpdfstring{\mr{समजूतदार} (samjūtdār)}{समजूतदार (samjūtdār)}]{\mr{समजूतदार} (samjūtdār) — understanding; mature}

\label{lesson:MR-C471-samjutdar}



\begin{quote}
Before the new one: say the Marathi for shrewd, then the Marathi for intelligent.
\end{quote}

\subsection*{You'll want to know: \mr{समजूतदार}}
\textbf{\mr{समजूतदार}} — \emph{samjūtdār} — \textquotedblleft{}understanding; mature\textquotedblright{}.

\begin{quote}
\mr{माझा} \mr{मोठा} \mr{भाऊ} \mr{खूप} \mr{समजूतदार} \mr{आहे}.
\end{quote}

\begin{quote}
\emph{mājhā moṭhā bhāū khūp samjūtdār āhe.}
\end{quote}

\begin{quote}
My big brother is very understanding.
\end{quote}

\mr{समजूत}, understanding, from \mr{समजणे}, + -\mr{दार}, having.

\begin{cousinweb}
\textbf{\mr{समजूतदार}} (\emph{samjūtdār}) is a Marathi derivation: \mr{समजूत}, \textquotedblleft{}understanding\textquotedblright{}, the noun of \mr{समजणे}, with \emph{-dār}, \textquotedblleft{}having\textquotedblright{}, an ending usually traced to Persian \emph{-dār}.
\end{cousinweb}

\begin{grammarlens}[title={What you've built}]
One more word for describing what you lived and felt.
\end{grammarlens}

\subsection*{Your turn}
\begin{itemize}
\raggedright
  \item \emph{Say it:} \emph{samjūtdār}
  \item \emph{Say it:} \emph{samjūtdār}, once more
  \item \emph{Say it:} \emph{mājhā moṭhā bhāū khūp samjūtdār āhe.} — My big brother is very understanding.
  \item \emph{Recall:} say the Marathi for shrewd, then the Marathi for intelligent, then say \emph{samjūtdār} again
\end{itemize}

\begin{tcolorbox}[breakable,colback=teal!4,colframe=teal!35!black,title={Before you move on}]
What does \mr{समजूतदार} mean? (\textquotedblleft{}Understanding; mature\textquotedblright{}.) Say it once more.
\end{tcolorbox}
\section[\texorpdfstring{\mr{अक्कल} (akkal)}{अक्कल (akkal)}]{\mr{अक्कल} (akkal) — common sense; wits}

\label{lesson:MR-C471-akkal}



\begin{quote}
Before the new one: say the Marathi for intelligent, then the Marathi for understanding.
\end{quote}

\subsection*{You'll want to know: \mr{अक्कल}}
\textbf{\mr{अक्कल}} — \emph{akkal} — \textquotedblleft{}common sense; wits\textquotedblright{}.

\begin{quote}
\mr{या} \mr{कामाला} \mr{थोडी} \mr{अक्कल} \mr{लागते}.
\end{quote}

\begin{quote}
\emph{yā kāmālā thoḍī akkal lāgte.}
\end{quote}

\begin{quote}
This job takes a little common sense.
\end{quote}

Feminine.

\mr{अक्कल} \mr{लढवणे}, to use one's wits; \mr{अक्कलहुशार}, quick-witted.

\begin{cousinweb}
\textbf{\mr{अक्कल}} (\emph{akkal}) is a loan from Persian \emph{aql}, \textquotedblleft{}reason\textquotedblright{}, usually traced to Arabic.
\end{cousinweb}

\begin{grammarlens}[title={What you've built}]
One more word for describing what you lived and felt.
\end{grammarlens}

\subsection*{Your turn}
\begin{itemize}
\raggedright
  \item \emph{Say it:} \emph{akkal}
  \item \emph{Say it:} \emph{akkal}, once more
  \item \emph{Say it:} \emph{yā kāmālā thoḍī akkal lāgte.} — This job takes a little common sense.
  \item \emph{Recall:} say the Marathi for intelligent, then the Marathi for understanding, then say \emph{akkal} again
\end{itemize}

\begin{tcolorbox}[breakable,colback=teal!4,colframe=teal!35!black,title={Before you move on}]
What does \mr{अक्कल} mean? (\textquotedblleft{}Common sense; wits\textquotedblright{}.) Say it once more.
\end{tcolorbox}
